%%%PAGE 297 rot=0 legibility=good summary=剪贴的印刷题(水平导轨上导体棒电流不变的电磁感应三问，带手写圈画与图上标注)；页面其余为空白，左缘露出另一剪贴纸条的残字
%%%BLOCK id=297-01 ch=12 sec=电磁感应·电路与能量 kind=problem alt=06 cont=none y=0.03-0.30
% 以下为剪贴的印刷题（题号被裁掉，仅剩一个“.”），其上有手写圈画；第1行末“长度”上方有一处手写小记号，第2行行首“小于”处有一道手写斜线
\begin{problem}
如图，水平\underline{面内有一光滑}金属导轨，其 $MN$、$PQ$ 边的电阻不计，$MP$ 边的电阻阻值 \fbox{$R=1.5\,\Omega$}，$MN$ 与 $MP$ 的\fbox{夹角为 $135^\circ$}，$PQ$ 与 $MP$ 垂直，$MP$ 边长度小于 $1\unit{m}$。将\fbox{质量 $m=2\unit{kg}$}、\fbox{电阻不计}的足够长直导体棒搁在导轨上，并与 $MP$ 平行。棒与 $MN$、$PQ$ 交点 $G$、$H$ 间的距离 \fbox{$L=4\unit{m}$}，空间存在垂直于导轨平面的匀强磁场，磁感应强度 \fbox{$B=0.5\unit{T}$}。在外力作用下，棒由 $GH$ 处以一定的初速度向左做直线运动，运动时回路中的电流强度始终与初始时的电流强度相等。

% 图上手写标注（已在图里）：v_3、2m、3m、4m、135° 等
%FIG p297_a 0.035 0.108 0.215 0.222
\notefig[0.35]{figures/p297_a.png}

\begin{enumerate}[label=(\arabic*)]
\item 若初速度 $v_1=3\unit{m/s}$，求棒在 $GH$ 处所受的安培力大小 $F_A$。
\item 若初速度 $v_2=1.5\unit{m/s}$，求棒向左移动距离 $2\unit{m}$ 到达 $EF$ 所需时间 \unclear{$\Delta t$}。
\item 在棒由 $GH$ 处向左移动 $2\unit{m}$ 到达 $EF$ 处的过程中，外力做功 \unclear{$W=7\unit{J}$}，求初速度 $v_3$。
\end{enumerate}
\end{problem}
%%%END
%%%BLOCK id=297-02 ch=99 sec=纸条残片 kind=misc alt=none cont=none y=0.57-0.85
页面左缘一条灰色纸条上露出的印刷字，自上而下：“.”、向、\unclear{愿}、\unclear{身}、夹、在、A.、C.（疑为第295页所贴印刷题各行行首字“向／感／射／夹／在／A.／C.”的残边，无其它内容）
%%%END
%%%PAGEEND 297
